\section{Results}
\label{sec:results}
\suppressfloats[t]

\subsection{Adaptive Validation Improves, but Is Not a Final Test}

Across the optimization history, the starting incumbent's validation F1 of
.5729 rises to .8091 for the final promoted candidate (+41.2\% relative).
Recorded tokens fall from 45.11M to 42.47M (5.9\%), provider cost from
\$7.788 to \$2.585 (66.8\%), and terminal failures from 23/72 to 0. These
endpoints look jointly favorable, but they were selected after repeated
validation access and an adaptive threshold change.

Figure~\ref{fig:optimization-trajectory} separates the reported lineage from
two exploratory continuations and labels every evaluated round's decision.
Figure~\ref{fig:final-surface-edits} gives the complete initial and final text
of the three surfaces whose edits survive in the final harness. The light-gray
branches are not additional optimizer seeds: they are abandoned continuations
in which round 3 or the original $\tau_{\mathrm{imp}}=0$ round-6 decision was
treated as promoted.

\begin{figure*}[t]
\centering
\begin{minipage}{.46\linewidth}
\centering
\begin{tikzpicture}[
  x=.64cm,y=12cm,
  accepted/.style={circle,draw=editgreen,fill=white,line width=1.15pt,minimum size=5.4pt,inner sep=0pt},
  rejected/.style={cross out,draw=fixgrey,line width=1.05pt,minimum size=5.4pt,inner sep=0pt},
  every node/.style={font=\scriptsize}
]
  \draw[black!18] (0,0) grid[xstep=1,ystep=.05] (10,.35);
  \draw[->,black!70] (0,0) -- (10.35,0) node[right] {Round};
  \draw[->,black!70] (0,0) -- (0,.39);
  \node[rotate=90] at (-.95,.18) {Validation F1};
  \foreach \x in {0,...,10} \draw (\x,0) -- +(0,-.005) node[below=2pt] {\x};
  \foreach \y/\lab in {0/.55,.05/.60,.10/.65,.15/.70,.20/.75,.25/.80,.30/.85,.35/.90}
    \draw (0,\y) -- +(-.04,0) node[left=2pt] {\lab};

  \fill[editfill,opacity=.55] (0,.36) rectangle (5.45,.385);
  \fill[blue!7] (5.55,.36) rectangle (10,.385);
  \node[font=\scriptsize\bfseries,text=editgreen] at (2.7,.3725)
    {rounds 1--5: $\tau_{\mathrm{imp}}=0$};
  \node[font=\scriptsize\bfseries,text=blue!55!black] at (7.8,.3725)
    {rounds 6--10: $\tau_{\mathrm{imp}}=.01$};
  \draw[densely dashed,red!65,line width=.9pt] (5.5,0) -- (5.5,.36);

  % Reported incumbent lineage.
  \draw[editgreen,line width=1.45pt]
    (0,.0229) -- (1,.0229) -- (1,.1849) -- (2,.1849) -- (2,.1855)
    -- (4,.1855) -- (4,.2413) -- (5,.2413) -- (5,.2823)
    -- (7,.2823) -- (7,.3005) -- (10,.3005) -- (10,.2591);

  \node[accepted] at (0,.0229) {};
  \node[accepted] at (1,.1849) {};
  \node[accepted] at (2,.1855) {};
  \node[accepted] at (4,.2413) {};
  \node[accepted] at (5,.2823) {};
  \node[accepted] at (7,.3005) {};
  \node[accepted] at (10,.2591) {};

  % Rejections use the compact baseline row convention of the Self-Harness
  % trajectory figure; their measured scores and context are in the side panel.
  \foreach \x/\lab in {3/R3,6/R6,8/R8,9/R9} {
    \node[rejected] at (\x,.010) {};
    \node[text=fixgrey,anchor=south] at (\x,.017) {\lab};
  }

  % Separately executed continuations, excluded from the reported lineage.
  \draw[fixgrey!62,densely dashed,line width=1.05pt]
    (2,.1855) -- (3,.2338) -- (4,.2338);
  \node[circle,draw=fixgrey!75,fill=white,line width=1pt,minimum size=5pt,inner sep=0pt] at (3,.2338) {};
  \node[rejected] at (4,.2338) {};

  \draw[fixgrey!62,densely dashed,line width=1.05pt]
    (5,.2823) -- (6,.2868) -- (7,.2442);
  \node[circle,draw=fixgrey!75,fill=white,line width=1pt,minimum size=5pt,inner sep=0pt] at (6,.2868) {};
  \node[rejected] at (7,.2442) {};

  % Marker legend, placed below the iteration axis as in Self-Harness.
  \node[accepted] at (1.30,-.060) {};
  \node[anchor=west] at (1.43,-.060) {accepted};
  \node[rejected] at (3.90,-.060) {};
  \node[anchor=west] at (4.03,-.060) {rejected};
  \draw[fixgrey!62,densely dashed,line width=1.05pt] (6.05,-.060) -- (6.45,-.060);
  \node[anchor=west] at (6.55,-.060) {abandoned branch};
\end{tikzpicture}
\end{minipage}%
\hfill
\begin{minipage}{.52\linewidth}
\centering\scriptsize
\setlength{\tabcolsep}{2.2pt}
\renewcommand{\arraystretch}{1.05}
\begin{tabularx}{\linewidth}{@{}lrrrrrX@{}}
\toprule
R & F1 & $\Delta$F1 & $\Delta$tok & Fail. & Dec. & What changed \\
\midrule
1  & .7349 & +.1620 & $-8.9\%$   & $23\to6$ & Prom. & emit user-ID pairs \\
2  & .7355 & +.0007 & $+11.7\%$  & $6\to1$  & Prom. & recover missing users \\
3  & .7838 & +.0483 & $-3.4\%$   & $1\to1$  & Rej.  & wrong-format regression \\
4  & .7913 & +.0558 & $-12.1\%$  & $1\to0$  & Prom. & confirm target categories \\
5  & .8323 & +.0410 & $+16.8\%$  & $0\to0$  & Prom. & correct detected errors \\
6  & .8368 & +.0045 & $-14.1\%$  & $0\to0$  & Rej.  & below revised $\tau_{\mathrm{imp}}=.01$ \\
7  & .8505 & +.0182 & $-33.8\%$  & $0\to0$  & Prom. & resolve count mismatch \\
8  & .8164 & $-$.0341 & $+56.5\%$  & $0\to1$  & Rej.  & format regression, +1 fail \\
9  & .6546 & $-$.1959 & $+665.7\%$ & $0\to7$  & Rej.  & token blowup, +7 fail \\
10 & .8091 & +.0128 & $-11.0\%$  & $0\to0$  & Prom. & rebuild after corrections \\
\bottomrule
\end{tabularx}
\end{minipage}
\caption{OOLONG-Pairs optimization trajectory (left) and the full per-round
record it summarizes (right). Changes are relative to that round's incumbent;
terminal failure counts are out of 72 attempts. Dark-green circles and the
step line denote accepted candidates on the reported incumbent lineage; gray
crosses in the row above the iteration axis denote rejected candidates.
Light-gray dashed traces show two abandoned exploratory continuations, found
by checking every sibling experiment directory in this lineage.}
\label{fig:optimization-trajectory}
\end{figure*}

% Diff-card primitives for fig:final-surface-edits: a full-width shaded row
% per logical line, "+" for appended text in that round's own color, plain
% (phantom-indented to match) for text retained verbatim from H0.
% Colorbox-around-parbox (not soul's \hl) so the shading spans the whole
% card width even when a line wraps. Round-to-text attribution verified
% against the real harness.json snapshots at each round (not the short
% behavioral_change summaries, which omit one of round 10's two insertions).
\newcommand{\diffsame}[1]{%
  \noindent\parbox{\dimexpr\linewidth-2\fboxsep\relax}{\ttfamily\small\phantom{+ }#1}\par\vspace{2pt}}
\newcommand{\diffadd}[2]{%
  \noindent\colorbox{#1}{\parbox{\dimexpr\linewidth-2\fboxsep\relax}{\ttfamily\small+\ #2}}\par\vspace{2pt}}
% Round-number badge for card titles: black text, background matching that
% round's own row fill color (not colored text) -- tight padding so it reads
% as a small tag rather than a full-size colorbox.
\newcommand{\rbadge}[2]{{\setlength{\fboxsep}{1.5pt}\colorbox{#1}{\textcolor{black}{#2}}}}

\begin{figure*}[p]
\centering
\setlength{\fboxsep}{4pt}
\begin{tcolorbox}[
  enhanced, colback=white, colframe=black!35, boxrule=0.5pt,
  arc=2pt, left=3pt, right=3pt, top=3pt, bottom=4pt,
  colbacktitle=black!8, coltitle=black, fonttitle=\bfseries\small,
  title={S2---decomposition instruction (promoted round \rbadge{r4fill}{4})}
]
\diffsame{Start by probing `context' to understand what you have, then decide how the task breaks into steps.}
\diffadd{r4fill}{Before selecting target categories, re-read the task question from `context' to identify which categories are required.}
\diffadd{r4fill}{Store the required categories in a variable and print them to confirm they match the task.}
\end{tcolorbox}

\vspace{4pt}
\begin{tcolorbox}[
  enhanced, colback=white, colframe=black!35, boxrule=0.5pt,
  arc=2pt, left=3pt, right=3pt, top=3pt, bottom=4pt,
  colbacktitle=black!8, coltitle=black, fonttitle=\bfseries\small,
  title={S3---execution instruction (cumulative promoted edits from rounds
    \rbadge{r1fill}{1}, \rbadge{r2fill}{2}, \rbadge{r5fill}{5}, and
    \rbadge{r10fill}{10})}
]
\diffsame{On each turn, run one block of code, print what you need to see, and build the result up in REPL variables.}
\diffadd{r1fill}{After classifying questions, generate the required user-ID pairs and output them in the format (user\_id\_1, user\_id\_2), one per line.}
\diffadd{r10fill}{When generating pairs from two user sets, skip pairs where the two user IDs are equal (a user cannot pair with itself).}
\diffadd{r2fill}{After classifying, verify that all users with entity or human being questions are captured. If any user appears to be missing, reclassify their questions.}
\diffadd{r5fill}{If you detect a classification error (e.g., a user's question count doesn't match the expected pattern), reclassify the affected questions before proceeding.}
\diffadd{r10fill}{After any reclassification, you MUST rebuild the user-to-category mapping and regenerate all pairs from scratch using the updated classifications. Do not reuse a previously computed user set or pair list after any classification change.}
\end{tcolorbox}

\vspace{4pt}
\begin{tcolorbox}[
  enhanced, colback=white, colframe=black!35, boxrule=0.5pt,
  arc=2pt, left=3pt, right=3pt, top=3pt, bottom=4pt,
  colbacktitle=black!8, coltitle=black, fonttitle=\bfseries\small,
  title={S4---verification instruction (promoted round \rbadge{r7fill}{7})}
]
\diffsame{Check your candidate answer before setting `answer["ready"] = True'.}
\diffadd{r7fill}{If you compute an expected count (e.g., of pairs, users, or classifications) and it does not match the actual count, do not submit until the discrepancy is resolved.}
\diffadd{r7fill}{Print the source of the discrepancy: which users or classifications differ between the expected and actual sets.}
\diffadd{r7fill}{Correct the underlying issue (e.g., reclassify affected questions, regenerate pairs) and re-verify before submitting.}
\end{tcolorbox}

\caption{Full retained field values for every surface changed in the frozen
round-10 harness, shown as a diff against the initial $H_0$ text: a plain line
is retained verbatim from $H_0$; a ``+'' line is new text, colored by the
round that added it (round number colored to match in each card's title).
S3 draws on four rounds, each shown in its own color; round 10 contributes
two separate insertions, one in the middle of the existing text and one at
the end, not a single trailing append. Every edit to these three surfaces was
a pure addition -- no span of the original text was altered or removed, so
no ``-'' lines appear anywhere. Unlike Figure~\ref{fig:optimization-trajectory},
this figure does not summarize the edits: it reproduces the full text from
the saved harness artifacts, diffed directly against the real per-round
harness snapshots (not the short per-candidate change summaries, which omit
one of round 10's two insertions).}
\label{fig:final-surface-edits}
\end{figure*}


Figure~\ref{fig:optimization-trajectory}'s side panel also demonstrates why
this run is not a Pareto-aware optimizer. Round 6 improves every displayed
objective (F1, tokens, failures) but is rejected by the scalar threshold,
while round 7 is promoted despite a 65.8\% provider-cost increase (its token
count falls 33.8\%). Applying Equation~\ref{eq:pareto-audit} to the available
round-level objectives, round 6 dominates its incumbent, whereas round 7 and its incumbent are
non-dominated because quality and tokens improve but cost worsens. The scalar
gate therefore produces \textbf{opposite decisions to the Pareto relation in
both rounds}. The appropriate claim is \emph{post hoc Pareto evaluation of a
constrained quality optimizer}.

\subsection{Pareto-Aware Frozen Length Audit}

\begin{table}[ht]
\centering\small
\setlength{\tabcolsep}{3.5pt}
\begin{tabular}{@{}lrrrrr@{}}
\toprule
Method & F1 & Tokens & Cost & Fail. & Instab. \\
\midrule
\hzero{}      & .6967 & 53.55M & \$2.788 & 6 & .2289 \\
\hzerostar{}  & .7248 & \textbf{26.72M} & \textbf{\$1.733} & 6 & .2048 \\
\shrlm{}      & \textbf{.7442} & 58.08M & \$2.915 & \textbf{2} & \textbf{.1630} \\
$\lambda$-RLM & .8472 &  1.05M & \$0.144 & 0 & .0321 \\
\bottomrule
\end{tabular}
\caption{Frozen 48-task OOLONG-Pairs test (96 runs/method). Tokens and cost are
totals; Fail. counts terminal failures. Instab. is the mean within-task absolute
difference between two attempts. $\lambda$-RLM is task-specific and its host
compute is unpriced. Bold marks the best generic harness on each objective.}
\label{tab:frozen48}
\end{table}

\begin{table}[ht]
\centering\scriptsize
\setlength{\tabcolsep}{3.2pt}
\renewcommand{\arraystretch}{1.02}
\begin{tabular}{@{}llrrrr@{}}
\toprule
Method & Metric & 8K & 16K & 32K & Macro / total \\
\midrule
\hzero{} & F1 & .7891 & .6481 & .6529 & .6967 \\
 & Tokens & 19.03M & 20.99M & 13.52M & 53.55M \\
 & Cost & \$0.816 & \$1.071 & \$0.900 & \$2.788 \\
\addlinespace[2pt]
\hzerostar{} & F1 & .6504 & \textbf{.7591} & \textbf{.7649} & .7248 \\
 & Tokens & \textbf{10.39M} & \textbf{8.81M} & \textbf{7.53M} & \textbf{26.72M} \\
 & Cost & \textbf{\$0.496} & \textbf{\$0.529} & \textbf{\$0.708} & \textbf{\$1.733} \\
\addlinespace[2pt]
\shrlm{} & F1 & \textbf{.8185} & .7142 & .7000 & \textbf{.7442} \\
 & Tokens & 18.07M & 22.40M & 17.61M & 58.08M \\
 & Cost & \$0.725 & \$0.931 & \$1.259 & \$2.915 \\
\addlinespace[2pt]
$\lambda$-RLM & F1 & .8829 & .8440 & .8148 & .8472 \\
 & Tokens & 0.129M & 0.328M & 0.592M & 1.05M \\
 & Cost & \$0.016 & \$0.050 & \$0.078 & \$0.144 \\
\bottomrule
\end{tabular}
\caption{Per-length frozen results for all four conditions. Each cell reports
mean F1, total recorded tokens, or total provider cost. Each length contains
16 tasks and 32 runs per method; the final column is macro F1 or the total over
all 96 runs. Bold marks the best generic harness; $\lambda$-RLM is specialized
and its host computation is unpriced.}
\label{tab:frozen48-length}
\end{table}

On the frozen test (Tables~\ref{tab:frozen48}--\ref{tab:frozen48-length}), \hzerostar{} strictly dominates
\hzero{}: it has higher F1, about half the tokens, 38\% lower cost, equal
failures, and lower instability. \shrlm{} raises observed F1 by .0194 over
\hzerostar{} and reduces failures, but uses 117\% more tokens and costs 68\%
more; neither dominates the other. The paired task-bootstrap interval for the
F1 difference is $[-.0521,.0935]$. By length, \shrlm{}--\hzerostar{} is +.1680
at 8K (95\% CI $[.0352,.3143]$), $-.0450$ at 16K
($[-.1590,.0561]$), and $-.0649$ at 32K ($[-.1586,.0278]$). Relative to
\hzerostar{}, \shrlm{} uses \textbf{74\%, 154\%, and 134\% more tokens} and
costs 46\%, 76\%, and 78\% more at 8K, 16K, and 32K, respectively. Thus the
aggregate gain comes from the 8K stratum and does not establish monotonically
improving length generalization or a quality--resource dominance relation.

$\lambda$-RLM exceeds \shrlm{} by .1031 F1 (95\% CI $[.0580,.1534]$), with
positive differences in all three length strata. It also records no terminal
failures and costs \$0.016/\$0.050/\$0.078 across 8K/16K/32K, compared with
\shrlm{}'s \$0.725/\$0.931/\$1.259. This is strong evidence that task-specific
algorithmic structure matters. It is not a clean systems-cost dominance claim,
because its deterministic host work is excluded from provider accounting.

\subsection{Long-Context and Cross-Task Audits}

The 8K/16K/32K division applies only to the stratified frozen length test.
The following evaluations each have one fixed regime: nominal 262K context for
OOLONG-Pairs and a full 3,507-document candidate corpus for OBLIQ-Bench Math.

\begin{table}[ht]
\centering\small
\setlength{\tabcolsep}{4pt}
\begin{tabular}{@{}lrrrr@{}}
\toprule
Method & 262K F1 & Tokens & Cost & Fail. \\
\midrule
\hzero{}      & .5696 & \textbf{80.13M} & \$10.962 & 7 \\
\hzerostar{}  & .5923 & 80.53M & \textbf{\$9.154}  & \textbf{1} \\
\shrlm{}      & \textbf{.6307} & 87.95M & \$10.384 & 2 \\
$\lambda$-RLM & .7465 & 22.12M & \$4.816  & 0 \\
\bottomrule
\end{tabular}
\caption{Frozen 40-task, nominal 262K OOLONG-Pairs test (80 runs/method).
Bold marks the best generic harness on each objective.}
\label{tab:long}
\end{table}

The 262K test repeats the frontier pattern (Table~\ref{tab:long}). \shrlm{}
improves over \hzerostar{} by .0383 F1, but the 95\% interval
$[-.0417,.1212]$ crosses zero and resource use is higher. $\lambda$-RLM's
.1158 advantage over \shrlm{} is resolved ($[.0737,.1657]$).

\begin{table}[ht]
\centering\scriptsize
\setlength{\tabcolsep}{3.5pt}
\begin{tabular}{@{}lrrrr@{}}
\toprule
Method & NDCG@10 & Tokens & Cost & Wrong fmt. \\
\midrule
\hzerostar{} & \textbf{.1040} & 70.18M & \$3.965 & \textbf{7} \\
\hzero{} & .0850 & 51.11M & \textbf{\$2.485} & 18 \\
\shrlm{} & .0517 & \textbf{50.54M} & \$2.534 & 29 \\
\bottomrule
\end{tabular}
\caption{Cross-task transfer to the frozen 40-query, full-corpus OBLIQ-Bench
Math split (80 runs/method). Tokens and provider cost are totals; Wrong fmt.
counts unparsable final answers. \shrlm{} receives no OBLIQ feedback. Bold
marks the best value on each objective.}
\label{tab:obliq-transfer}
\end{table}

Transfer is negative (Table~\ref{tab:obliq-transfer}). Relative to
\hzerostar{}, the frozen round-10 harness decreases mean NDCG@10 by 50.2\%
($-.0522$, paired 95\% CI $[-.0914,-.0175]$) and raises wrong-format failures
from 7/80 to 29/80. This loss is resolved despite \shrlm{} using 28.0\% fewer
tokens and 36.1\% lower provider cost. \hzerostar{} also has higher observed
NDCG@10 than \hzero{} by .0190, although its interval
$[-.0051,.0447]$ crosses zero. Because no OBLIQ feedback enters the search,
the result provides a concrete counterexample to broad cross-task robustness
claims for the learned harness.

\section{Discussion and Conclusion}

The optimization run is successful as an engineering search: it discovers a
harness with much better adaptive validation behavior and fewer formatting
failures. It is not evidence that the final harness is unconditionally better.
Frozen evaluation instead identifies a quality--efficiency frontier, unresolved
quality differences between generic harnesses, a decisive advantage for a
task-specific algorithm, and negative transfer. Reproducible AI reporting should
therefore distinguish three objects: the adaptive optimization record, an
untouched frozen test, and the multi-objective deployment choice. Publishing all
attempts and failures changes the scientific conclusion from ``automatic harness
improvement'' to the narrower and better-supported finding: scalar-guided
self-harnessing can find useful trade-offs, but robust improvement requires an
external Pareto audit.
